%! Project: Design and Conduct a Survey — Unit 1, Lesson 1.8 — Slides (Beamer)
% PILOT SHAPE (2026-09-06): modeled on AP Statistics lesson 1.4 minus its AP Practice page.
% GRADUAL RELEASE deck: title → targets → warm-up → hook (unresolved) → I DO divider → two or
% three frames per notes row → WE DO (the class's own survey — a LIVE surface, un-answered) →
% spoken debrief with the cold check → close & START THE HOMEWORK, whose item 6 is the project
% deliverable. No group activity, no practice launch, no exit ticket. There is no spoiler rule:
% the notes frames ARE the direct instruction.
% \CourseName is NOT defined in beamer — the course name is written literally.
\documentclass[aspectratio=169, 11pt]{beamer}
\usepackage{saar-beamer}   % provides \ceruleanheader{} and \sectionlabel[color]{}

\newcommand{\redink}[1]{{\color{redacc}\bfseries #1}}

\begin{document}

% TITLE SLIDE
{
\setbeamercolor{background canvas}{bg=cerulean}
\begin{frame}[plain]
  \vspace{2.8cm}\hspace{0.55cm}%
  \begin{minipage}{0.9\textwidth}
    {\color{goldacc}\bfseries\small LESSON 1.8}\\[0.5cm]
    {\color{white}\bfseries\LARGE Project: Design and Conduct a Survey}\\[1.2cm]
    {\color{frostmid}\small Statistical Analysis \& Algebraic Reasoning~$\cdot$~Unit 1}
  \end{minipage}
\end{frame}
}

% ── LEARNING TARGETS ─────────────────────────────────────────────────────────
\begin{frame}[t, plain]
  \ceruleanheader{Today's targets}
  \sectionlabel{Lesson 1.8 $\cdot$ the unit, run forwards}
  By the end of today, I can\ldots
  \begin{itemize}\setlength{\itemsep}{5pt}
    \item plan a survey through the four stages of the \textbf{statistical cycle}, and say what
          each stage hands in.
    \item write a \textbf{survey instrument} whose \textbf{closed-response items} ask one thing,
          in neutral words, with choices that cover everyone.
    \item turn a tally into \textbf{relative frequencies} and read them off a \textbf{bar graph}.
    \item report a conclusion about \emph{the people my design actually covers}.
  \end{itemize}
  \begin{block}{How today works}
    Warm-up $\to$ \textbf{notes together} $\to$ \textbf{we build and run our own survey} $\to$
    debrief $\to$ \textbf{you start the homework here, alone}. Item 6 of the homework is your
    project report.
  \end{block}
\end{frame}

% ── WARM-UP ──────────────────────────────────────────────────────────────────
\begin{frame}[t, plain]
  \ceruleanheader{Warm-Up --- 5 minutes, silent}
  \sectionlabel{the Student Council's finished survey}
  \vspace{0.25cm}
  \begin{center}
    \begin{minipage}{0.88\textwidth}
      A \$3{,}000 grant. The council drew a \textbf{stratified sample of 60} of Riverbend's
      \textbf{900} students: club funding \textbf{21}, field trips 18, gym equipment 15, longer
      lunch 6.
      \begin{enumerate}\setlength{\itemsep}{6pt}
        \item What percent of the 60 chose club funding?
        \item Predict how many of all 900 would choose it.
        \item Two groups are hiding in that paragraph --- how many in each? And how many choices
              did each student pick?
      \end{enumerate}
    \end{minipage}
  \end{center}
  \vspace{0.3cm}
  \begin{center}
    {\color{cerulean}\bfseries 60 asked. 900 reported on. That line stays on the board all period.}
  \end{center}
\end{frame}

% ── HOOK — leave it UNRESOLVED; problem 13 (row 4, The Limit) settles it ─────
{
\setbeamercolor{background canvas}{bg=cerulean}
\begin{frame}[plain]
  \vspace{1.1cm}
  \begin{center}
    {\color{frostmid}\bfseries\small LAST SPRING, A CLASS OF 25 IN THIS ROOM ASKED ITSELF THE SAME QUESTION.}\\[0.7cm]
    {\color{white}\bfseries\Large The class got 36\%. \quad The council got 35\%.}\\[0.6cm]
    {\color{white}\small ``We matched the council, so our class can speak for all 900 students too.''}\\[0.8cm]
    {\color{goldacc}\bfseries\large Is that student right?}\\[0.4cm]
    {\color{frostmid}\small Vote: agree or disagree. We settle this at problem 13.}
  \end{center}
\end{frame}
}

% ── I DO — direct instruction begins ─────────────────────────────────────────
{
\setbeamercolor{background canvas}{bg=cerulean}
\begin{frame}[plain]
  \vspace{1.8cm}\hspace{0.8cm}%
  \begin{minipage}{0.86\textwidth}
    {\color{goldacc}\bfseries\small GUIDED NOTES \ \textbullet\ \ I DO}\\[0.7cm]
    {\color{white}\bfseries\Large Open to your Guided Notes page. We work the table row by row
    --- fill each term as we name it, not before.}\\[0.9cm]
    {\color{frostmid}\small Five terms today. The last row is \emph{our} survey --- we build it,
    we answer it, we tally it.}
  \end{minipage}
\end{frame}
}

% ── 1. THE CYCLE ─────────────────────────────────────────────────────────────
\begin{frame}[t, plain]
  \ceruleanheader{Your project is the cycle, run forwards}
  \sectionlabel{notes row 1 $\cdot$ The Cycle $\cdot$ problems 1--4}
  \vspace{2pt}
  \begin{center}
    \small
    \renewcommand{\arraystretch}{1.3}
    \begin{tabular}{@{} l l @{}}
    \hline
    \textbf{Stage} & \textbf{What you hand in} \\
    \hline
    1. Formulate        & one statistical question, and the population it is about \\
    2. Collect          & a sampling plan, and the finished instrument \\
    3. Represent        & a frequency table, and a bar graph \\
    4. Analyze \& report & a written report with its limits named \\
    \hline
    \end{tabular}
  \end{center}
  \vspace{8pt}
  \begin{itemize}\setlength{\itemsep}{4pt}
    \item A survey is an \textbf{observational study} --- nobody is assigned anything.
    \item \textbf{Problem 4:} our class will ask only this room. Who chose that sample ---
          chance, or the schedule? \redink{A convenience sample.}
  \end{itemize}
\end{frame}

% ── 2a. THE INSTRUMENT — the four rules ──────────────────────────────────────
\begin{frame}[t, plain]
  \ceruleanheader{The instrument is the wording, and the wording is the data}
  \sectionlabel{notes row 2 $\cdot$ The Instrument $\cdot$ problems 5--6}
  \vspace{2pt}
  \begin{itemize}\setlength{\itemsep}{5pt}
    \item \textbf{Rule 1 --- one idea per item.} Two ideas cannot be answered by one person.
    \item \textbf{Rule 2 --- neutral wording.} The item must not tell the person what you want.
    \item \textbf{Rule 3 --- choices cover everyone and do not overlap.} Then, and only then, the
          percents add to 100\%.
    \item \textbf{Rule 4 --- short.} A person answers what they finish reading.
  \end{itemize}
  \vspace{6pt}
  \begin{block}{Why this row is the one that cannot be fixed later}
    A biased instrument is wrong \emph{in a direction}. \redink{No arithmetic done afterwards
    pulls it back} --- you are dividing the wrong counts by the right total.
  \end{block}
\end{frame}

% ── 2b. THE INSTRUMENT — repair the three broken items ───────────────────────
\begin{frame}[t, plain]
  \ceruleanheader{Three broken items --- you met all three in Lesson 1.4}
  \sectionlabel{notes row 2 $\cdot$ the fill table $\cdot$ problems 5--6}
  \vspace{2pt}
  \begin{center}
    \small
    \renewcommand{\arraystretch}{1.35}
    \begin{tabular}{@{} p{7.2cm} l l @{}}
    \hline
    \textbf{The item as written} & \textbf{Rule} & \textbf{Its 1.4 name} \\
    \hline
    ``Don't you agree that our underfunded clubs deserve the grant?''
      & \redink{rule 2} & \redink{leading} \\
    ``Should the council fund clubs and field trips?''
      & \redink{rule 1} & \redink{double-barrelled} \\
    ``How old are you: 14--16, 16--18, or 18+?''
      & \redink{rule 3} & \redink{overlapping} \\
    \hline
    \end{tabular}
  \end{center}
  \vspace{8pt}
  \begin{center}
    {\color{cerulean}\bfseries Where does a 16-year-old check? Now add the percents.}
  \end{center}
\end{frame}

% ── 3. RELATIVE FREQUENCY ────────────────────────────────────────────────────
\begin{frame}[t, plain]
  \ceruleanheader{Counts $\to$ percents $\to$ picture}
  \sectionlabel{notes row 3 $\cdot$ Relative Frequency $\cdot$ problems 7--11}
  \vspace{2pt}
  \begin{center}
    \small
    \renewcommand{\arraystretch}{1.25}
    \begin{tabular}{@{} l c c @{}}
    \hline
    \textbf{Choice} & \textbf{Frequency} & \textbf{Relative frequency} \\
    \hline
    Club funding  & 21 & $21/60 =$ \redink{35\%} \\
    Field trips   & 18 & $18/60 =$ \redink{30\%} \\
    Gym equipment & 15 & $15/60 =$ \redink{25\%} \\
    Longer lunch  & \phantom{0}6 & $6/60 =$ \redink{10\%} \\
    \hline
    \textbf{Total} & 60 & \redink{100\%} \\
    \hline
    \end{tabular}
  \end{center}
  \vspace{6pt}
  \begin{itemize}\setlength{\itemsep}{4pt}
    \item ``Choice'' is \textbf{categorical}, so the display is a \textbf{bar graph} --- bars
          separated, \textbf{height $=$ relative frequency}.
    \item \textbf{Problem 10:} 42 of 120 is \emph{also} 35\%. The bar does not move. \redink{A
          height is a share, not a count.}
  \end{itemize}
\end{frame}

% ── 3b. THE TRAP ─────────────────────────────────────────────────────────────
\begin{frame}[t, plain]
  \ceruleanheader{The arithmetic is right. The claim is not.}
  \sectionlabel{notes row 3 $\cdot$ problem 9 $\cdot$ the trap --- vote before you write}
  \vspace{4pt}
  \begin{center}
    \begin{minipage}{0.86\textwidth}
      Last spring's class of \textbf{25}: \textbf{9} chose club funding.
      \[ \frac{9}{25} = 0.36 = 36\% \qquad\qquad 0.36 \times 900 = 324 \]
    \end{minipage}
  \end{center}
  \vspace{4pt}
  \begin{itemize}\setlength{\itemsep}{5pt}
    \item Is either computation wrong? \redink{No. Both are correct.}
    \item May the class report that 324 about the school? \redink{No.}
    \item Who does 324 describe? \redink{Nobody.} The schedule chose the class, not chance.
  \end{itemize}
\end{frame}

% ── 4. THE CRUX ──────────────────────────────────────────────────────────────
\begin{frame}[t, plain]
  \ceruleanheader{Near-agreement is luck, not evidence}
  \sectionlabel[redacc]{notes row 4 $\cdot$ The Limit $\cdot$ problems 12--15 $\cdot$ the whole point of today}
  \vspace{2pt}
  \begin{columns}[t]
    \begin{column}{0.47\textwidth}
      \begin{block}{The council --- 60 of 900}
        Stratified by grade. \textbf{Chance} chose them. \\[3pt]
        $21/60 = 35\%$ \\[3pt]
        May describe: \textbf{all 900 students.}
      \end{block}
    \end{column}
    \begin{column}{0.47\textwidth}
      \begin{block}{The class --- 25 of 900}
        One room. The \textbf{schedule} chose them. \\[3pt]
        $9/25 = 36\%$ \\[3pt]
        May describe: \redink{this class only.}
      \end{block}
    \end{column}
  \end{columns}
  \vspace{8pt}
  \begin{center}
    {\large \textbf{Problem 13 --- vote again: agree or disagree?}}\\[5pt]
    {\small\color{slate} Chance never chose the class, so it could have missed by twenty points
    --- and nothing in its method would have told anyone. \textbf{The match happened; it was not
    earned.}}
  \end{center}
\end{frame}

% ── WE DO — a LIVE surface: the class's own survey, un-answered ──────────────
\begin{frame}[t, plain]
  \ceruleanheader{Guided Practice: Our Class Survey}
  \sectionlabel[goldacc]{We do $\cdot$ problems 16--19 $\cdot$ you hold the pen}
  \vspace{2pt}
  \begin{itemize}\setlength{\itemsep}{4pt}
    \item[] \textbf{3 min:} agree on \emph{one} item and four choices. Check it against all four
            rules before anyone writes it down.
    \item[] \textbf{4 min:} ask every student in this room. Tally. Divide each count by the total.
    \item[16.] Population, sample, $n$ --- and who chose the sample?
    \item[17.] Our largest category: show the division, then write our sentence 4.
    \item[18.] The council's 60 and us: whose result may be reported about all 900? If ours lands
          near theirs, what does that prove?
    \item[19.] Which limit would you fix first --- and what would fixing it cost?
  \end{itemize}
  \begin{block}{The question behind the question}
    If you said the match proves something in 18, you are still at the hook vote.
  \end{block}
\end{frame}

% ── DEBRIEF ──────────────────────────────────────────────────────────────────
\begin{frame}[t, plain]
  \ceruleanheader{Debrief: whom may this report describe?}
  \sectionlabel{10 minutes $\cdot$ pens down, papers up --- we say these aloud}
  \vspace{2pt}
  \begin{enumerate}\setlength{\itemsep}{4pt}
    \item The two panels --- say the sentence: near-agreement is luck, not evidence.
    \item ``$0.36 \times 900 = 324$.'' Is it wrong? Then who does 324 describe?
    \item ``Don't you agree\ldots'' --- which rule, which 1.4 name, and why no later arithmetic
          repairs it.
  \end{enumerate}
  \vspace{4pt}
  \begin{block}{Cold check --- hands down, thirty seconds}
    A homeroom of \textbf{30} at Bayside Middle School surveyed itself about spirit week:
    decades 12, color wars 9, movie characters 6, sports teams 3. It told the whole school of
    \textbf{750}: ``Bayside students want a decades theme.'' The arithmetic checks out ---
    $12/30 = 40\%$ and $0.40 \times 750 = 300$. \textbf{Is the report honest? What would you
    change?}
  \end{block}
\end{frame}

% ── YOU DO — the homework is the individual practice, started in class ───────
\begin{frame}[t, plain]
  \ceruleanheader{You do: start the homework}
  \sectionlabel[goldacc]{Last 10 minutes $\cdot$ on your own $\cdot$ finish it at home}
  \begin{itemize}\setlength{\itemsep}{4pt}
    \item \textbf{Millbrook Public Library} --- 1{,}200 card holders. A phone \textbf{SRS of 80},
          and a front-desk \textbf{clipboard of 50}.
    \item Relative frequencies; read the bar graph; scale one of the two percents to all 1{,}200
          --- and say which one you may not scale.
    \item A clipboard at 26\% beside the survey's 25\%. Does the near-match earn it the library?
    \item \textbf{Item 6 is your project report} --- four sentences on \emph{our} tally. Take your
          notes page home; you need the numbers.
  \end{itemize}
  \begin{block}{Silent and alone --- I am circulating. Packet pages, scored out of 12, due at the start of the first class after two study halls.}
    Start with \textbf{1, 3, and 4}. Sentence 4 of item 6 is the one I am grading hardest ---
    its subject has to be the people we actually asked.
  \end{block}
\end{frame}

% ── CLOSE ────────────────────────────────────────────────────────────────────
{
\setbeamercolor{background canvas}{bg=cerulean}
\begin{frame}[plain]
  \vspace{1.5cm}\hspace{0.8cm}%
  \begin{minipage}{0.86\textwidth}
    {\color{goldacc}\bfseries\small BEFORE YOU GO}\\[0.7cm]
    {\color{white}\bfseries\Large You walked in able to take somebody else's study apart. You
    walk out having built one --- and knowing that the last sentence of your report is the part
    your design has to pay for.}\\[0.9cm]
    {\color{frostmid}\small Unit 1 ends here. Next: the unit review, then the test. Then Unit 2
    --- one \emph{quantitative} variable, where the bar graph gives way to a dot plot and the
    bars start to touch.}
  \end{minipage}
\end{frame}
}

\end{document}
